\documentclass[11pt]{article}

%============================ PREAMBLE ============================
\usepackage[T1]{fontenc}
\usepackage{lmodern}
\usepackage[a4paper,margin=1in]{geometry}
\usepackage{amsmath,amssymb}
\usepackage{booktabs}
\usepackage{longtable}
\usepackage{array}
\usepackage{xcolor}
\usepackage{listings}
\usepackage{caption}
\usepackage{enumitem}
\usepackage{microtype}
\usepackage{fancyhdr}
\usepackage{titlesec}
\usepackage[strings]{underscore}
\usepackage[hidelinks]{hyperref}
\hypersetup{colorlinks=true, linkcolor=black!70!blue, urlcolor=blue!60!black,
  pdftitle={Sleep, Activity, GERD and Oesophagitis (All of Us R9): Analysis Specification},
  pdfauthor={Functional GI Disorders and Physiological Measures Project}}

\DeclareMathOperator{\logit}{logit}

\definecolor{codebg}{RGB}{247,247,249}
\definecolor{codekw}{RGB}{0,90,160}
\definecolor{codestr}{RGB}{150,60,20}
\definecolor{codecom}{RGB}{90,120,90}
\definecolor{rulegray}{RGB}{120,120,120}
\definecolor{headblue}{RGB}{20,60,110}

\titleformat{\section}{\normalfont\large\bfseries\color{headblue}}{\thesection}{0.6em}{}
\titleformat{\subsection}{\normalfont\normalsize\bfseries\color{headblue!85!black}}{\thesubsection}{0.5em}{}
\titlespacing*{\section}{0pt}{1.4ex plus .2ex}{0.8ex}
\titlespacing*{\subsection}{0pt}{1.1ex plus .2ex}{0.6ex}

\pagestyle{fancy}\fancyhf{}
\renewcommand{\headrulewidth}{0.4pt}
\lhead{\footnotesize\color{rulegray}Sleep, Activity \& GERD/Oesophagitis --- All of Us R9}
\rhead{\footnotesize\color{rulegray}Analysis Specification}
\cfoot{\thepage}

\lstdefinelanguage{RR}{
  keywords={library,function,if,else,for,in,TRUE,FALSE,NA,NULL,return,source,readRDS,saveRDS,
    mutate,filter,select,group_by,summarise,left_join,inner_join,full_join,factor,case_when,
    glm,binomial,ntile,across,replace_na,stepAIC,as.Date,grepl,intersect,stopifnot},
  sensitive=true, morecomment=[l]{\#}, morestring=[b]", morestring=[b]'}
\lstset{language=RR, backgroundcolor=\color{codebg}, basicstyle=\ttfamily\scriptsize,
  keywordstyle=\color{codekw}\bfseries, stringstyle=\color{codestr},
  commentstyle=\color{codecom}\itshape, frame=single, rulecolor=\color{rulegray!50},
  breaklines=true, showstringspaces=false, columns=fullflexible, keepspaces=true,
  tabsize=2, captionpos=b, xleftmargin=0.8em,
  literate={>=}{{$\geq$}}1 {<=}{{$\leq$}}1}

\newcommand{\code}[1]{\texttt{\small #1}}

\title{\vspace{-1.3cm}\textbf{\color{headblue}Sleep, Physical Activity, and Acid-Related Upper-GI Disease}\\[2pt]
\textbf{\color{headblue}GERD and Oesophagitis in the \textit{All of Us} Research Program (v9)}\\[5pt]
\large Analysis Specification, Software Architecture, and Verification}
\author{Functional GI Disorders \& Physiological Measures Project}
\date{July 28, 2026}

\begin{document}
\maketitle
\thispagestyle{fancy}

\vspace{-0.6cm}
\begin{abstract}
\noindent This report specifies a retrospective cross-sectional analysis relating Fitbit-recorded
\emph{sleep} and \emph{physical-activity} metrics to two nested electronic-health-record (EHR)
outcomes --- broad gastroesophageal reflux disease (\code{has_gerd}) and its erosive
\emph{oesophagitis} subset (\code{has_esophagitis}) --- in the \emph{All of Us} Controlled Tier
CDR v9. The design is a $2\times3$ matrix: two outcomes crossed with three exposure sets (sleep,
activity, and a combined set), giving six analysis configurations executed by one shared statistical
engine. Seven R scripts implement the work without modifying any existing notebook. The analysis code
has been written against the \emph{verified} R9 schemas and validated end-to-end on schema-exact
data; this document records the phenotype definitions, cohort logic, covariate strategy, statistical
models, software architecture, and the verification results, including five defects found and fixed
during schema alignment. Outputs are emitted in the exact table and figure format required by the
target manuscript.
\end{abstract}

\vspace{-0.2cm}
{\small\noindent\textbf{Code.}
\code{GERD Outcome Data Upload R9.R},\;
\code{GERD Data Prep R9.R},\;
\code{GERD Analysis Helpers R9.R},\;
\code{Sleep and GERD Analysis File R9.R},\;
\code{Activity and GERD Analysis File R9.R},\;
\code{Activity Sleep and GERD Analysis File R9.R},\;
\code{GERD Pharmacologics Data Upload.R}.}

\vspace{0.2cm}
{\footnotesize\tableofcontents}
\newpage

%======================================================================
\section{Overview and Research Question}
%======================================================================
Gastroesophageal reflux disease (GERD) sits at the intersection of two disease constructs: an
acid-related structural disorder producing objective mucosal injury, and a symptom-defined
presentation sharing mechanisms with the disorders of gut--brain interaction (visceral
hypersensitivity, altered central pain processing, behavioural modulation). This dual character
motivates analysing \emph{two nested phenotypes} in parallel: broad, largely symptom-coded GERD, and
the erosive oesophagitis subset with documented mucosal damage.

The question is directional: \emph{sleep and physical activity} (exposures) $\rightarrow$ \emph{GERD
and oesophagitis} (outcomes). The prior hypothesis, consistent with the project's IBS, functional
dyspepsia and constipation studies, is that greater sleep quantity/quality and higher activity are
associated with \emph{lower} adjusted odds of both outcomes, attenuating but persisting after
adjustment.

\subsection{Candidate mechanisms}
These motivate both the covariate set and the directional predictions:
\begin{itemize}[leftmargin=1.3em,itemsep=0pt]
  \item \textbf{Acid clearance and supine time.} Swallowing and salivary bicarbonate are suppressed
  during sleep; fragmented or short sleep extends low-clearance recumbency and prolongs mucosal acid
  contact --- a pathway that should matter \emph{more} for the erosive phenotype.
  \item \textbf{Intra-abdominal pressure.} Higher BMI raises the gastro-oesophageal pressure gradient
  and promotes hiatal herniation, while also shortening sleep and reducing activity. This is the
  central confounder.
  \item \textbf{Lower-oesophageal-sphincter tone.} Nicotine, alcohol and sympathetic arousal increase
  transient sphincter relaxations.
  \item \textbf{Autonomic regulation.} Habitual activity raises vagal tone, which governs gastric
  motility and sphincter function; the corrected maximum-heart-rate exposure proxies this axis.
  \item \textbf{Visceral hypersensitivity.} Sleep deprivation lowers visceral pain thresholds, so
  identical acid exposure yields more symptoms --- acting preferentially on the \emph{symptom-coded}
  outcome. Contrasting the two outcomes therefore separates perceptual from structural pathways.
  \item \textbf{Reverse causation.} Nocturnal reflux itself fragments sleep. Adjustment cannot fix
  this; only design can, which is why a temporal outcome definition is the primary phenotype.
\end{itemize}

\subsection{Evidence gap}
Prior work measures sleep and activity largely by self-report and ascertains reflux by symptom
questionnaire in modest samples. Few datasets link high-resolution wearable streams to longitudinal
EHR at national scale. Notably, Hu et al.\ (2024) found subjective sleep quality differed sharply
across the GERD spectrum while objective measures did not --- a result that maps directly onto this
study's two-outcome design. This analysis contributes objectively measured exposures, an EHR-defined
outcome pair separating symptom-coded from erosive disease, and an unusually complete covariate set.

%======================================================================
\section{Study Design}
%======================================================================
\begin{table}[h]\centering\footnotesize
\begin{tabular}{@{}lp{10.6cm}@{}}
\toprule
Element & Specification \\ \midrule
Data source & \emph{All of Us} Controlled Tier CDR v9 (\code{fc-aou-cdr-prod-ct.C2024Q3R9}) \\
Design & Retrospective cross-sectional; wearable exposures, EHR outcomes \\
Population & Fitbit-sharers, sharing EHR, age $\geq 18$ at first Fitbit record \\
Exposures & Cohort quartiles (Q1 reference) of 8 sleep and 7 activity metrics \\
Outcomes & \code{has_gerd} and \code{has_esophagitis} (nested), $\geq 2$ records each \\
Primary rule & $\geq 2$ records with the \emph{first} $\geq 180$ days after the first Fitbit record \\
Estimator & Logistic regression; $e^{\beta_j}$ = adjusted OR for quartile $j$ vs Q1 \\
Multiplicity & Bonferroni (matching the antecedent IBS paper) \\
Configurations & Six: 2 outcomes $\times$ 3 exposure sets, one shared engine \\
\bottomrule
\end{tabular}
\caption{Design at a glance.}
\end{table}

\begin{table}[h]\centering\footnotesize
\begin{tabular}{@{}c p{2.6cm} p{2.9cm} p{2.2cm} p{3.4cm}@{}}
\toprule
ID & Outcome & Exposure set & Cohort & Analysis file \\ \midrule
C1 & Broad GERD   & Sleep (8) & $\geq 180$ nights & Sleep and GERD \\
C2 & Oesophagitis & Sleep (8) & $\geq 180$ nights & Sleep and GERD \\
C3 & Broad GERD   & Activity (7) & $\geq 30$ days & Activity and GERD \\
C4 & Oesophagitis & Activity (7) & $\geq 30$ days & Activity and GERD \\
C5 & Broad GERD   & Combined (15 + pairs) & Intersection & Activity Sleep and GERD \\
C6 & Oesophagitis & Combined (15 + pairs) & Intersection & Activity Sleep and GERD \\
\bottomrule
\end{tabular}
\caption{The six configurations. Configurations sharing a row-pair are produced by the same file,
because the driver loops both outcomes over one modelling frame.}
\label{tab:six}
\end{table}

%======================================================================
\section{Data Sources and Verified Schemas}
%======================================================================
Every column name and type below was verified by loading the dataset schema templates directly, not
inferred. This matters: several assumptions that looked reasonable turned out to be wrong
(Section~\ref{sec:verify}).

\subsection{Raw tables}
\begin{table}[h]\centering\footnotesize
\begin{tabular}{@{}p{5.0cm} r r p{5.4cm}@{}}
\toprule
Table & Rows & RAM & Columns used \\ \midrule
\code{R9_fitbit_sleep_daily_summary_df} & 31.8M & 2.4 GB & \code{sleep_date} (Date), \code{is_main_sleep}, the 7 \code{minute_*} \\
\code{R9_fitbit_activity_df} & 44.2M & 2.4 GB & \code{date} (Date), \code{*_active_minutes}, \code{sedentary_minutes} \\
\code{R9_fitbit_intraday_steps_df} & 34.2M & 1.0 GB & \code{date}, \code{wear_hours}, \code{total_steps} \\
\code{R9_drug_df} & 14.5M & 3.1 GB & \code{drug_class_concept_id}, \code{drug_exposure_start_datetime} \\
\code{R9_survey_df} & 9.7M & 0.7 GB & \code{question_concept_id}, \code{answer_concept_id} \\
\code{R9_condition_df} & 1.6M & 274 MB & \code{condition_concept_id}, \code{condition_start_datetime} \\
\code{R9_measurement_df} & 50,942 & 14 MB & \code{measurement_concept_id}, \code{value_as_number} \\
\code{R9_person_df} & 59,018 & 5 MB & \code{date_of_birth} (character), race, ethnicity, sex \\
\bottomrule
\end{tabular}
\caption{Raw inputs. Date-like EHR fields are stored as \emph{character} and require \code{as.Date()};
the Fitbit date fields are already \code{Date}.}
\end{table}

\subsection{Pre-built covariate files}
The project already contains outcome-agnostic covariate objects produced by the published IBS and
constipation pipelines. These are reused rather than recomputed (Section~\ref{sec:covstrategy}).

\begin{table}[h]\centering\footnotesize
\begin{tabular}{@{}p{5.4cm} r p{6.2cm}@{}}
\toprule
File & Rows & Contents \\ \midrule
\code{valid_population_sleep} & 19,995 & \textbf{the published cohort}: eligibility + age \\
\code{sleep_summary_filtered} & 23,812 & \code{n_valid_nights} + the 8 \code{avg_*} sleep metrics \\
\code{bmi_covariates_sleep_df} & 46,753 & \code{median_bmi}, \code{first_fitbit_sleep_date} \\
\code{cci_covariates_sleep_df} & 11,652 & \code{cci_score}, \code{cci_cat} \\
\code{medication_flags_sleep} & 11,105 & the 7 \code{on_*} flags \\
\code{comorbidity_status_sleep_ibs_df} & 8,079 & 9 \code{has_*} (no \code{has_pud}) \\
\code{R9_steps_summary_filtered} & 50,026 & \code{avg_daily_steps}, \code{n_valid_days} \\
\code{R9_activity_zone_summary} & 55,552 & zone minutes incl.\ \code{avg_total_active_min} \\
\code{R9_comorbidity_status_df} & 8,626 & 9 \code{has_*} \emph{plus} \code{has_pud} \\
\code{R9_pud_status} & \textbf{137} & \code{has_pud} --- very rare; see the sparsity guard \\
\bottomrule
\end{tabular}
\caption{Pre-built covariate inputs (sleep-anchored above, activity-anchored below the midline).}
\end{table}

\subsection{Join hazards}
\label{sec:joins}
Four schema features silently corrupt naive joins and are handled explicitly in the code:
\begin{enumerate}[leftmargin=1.4em,itemsep=1pt]
  \item \code{R9_avg_wear_hours_df} \emph{and} \code{R9_steps_summary_filtered} both carry
  \code{n_valid_days}; joining them yields \code{n_valid_days.x/.y} and breaks every later reference.
  Only \code{person_id, avg_daily_wear_hours} is taken.
  \item \code{alcohol_summary_df} carries \code{answer_concept_id.x} and \code{.y} from upstream
  joins; only \code{alcohol_likert_final} is taken.
  \item \code{education_df}/\code{income_df} carry \code{answer_concept_id}. Categories are mapped
  from \emph{concept IDs}, not the free-text \code{*_response}, whose values did not match the IBS
  script's string matching and silently collapsed everyone to ``Missing''.
  \item Factor levels use \textbf{en-dashes} (\code{age_cat} ``30--44''; \code{cci_cat} ``1--2'').
  Dash style is normalised before use.
\end{enumerate}

\noindent GERD is \emph{not} present in \code{R9_condition_df}, so a dedicated outcome pull is required.

%======================================================================
\section{Outcome Ascertainment}
%======================================================================
Both outcomes use the project's $\geq 2$-record specificity rule:
\begin{align*}
\texttt{has\_gerd} &= \mathbb{1}\{\geq 2 \text{ GERD condition records}\},\\
\texttt{has\_esophagitis} &= \mathbb{1}\{\geq 2 \text{ GERD records named ``esophagitis''}\}.
\end{align*}
Because every oesophagitis record is also a GERD record, \code{has_esophagitis} is a strict subset of
\code{has_gerd}; both are analysed as separate binary outcomes on the full cohort.

\subsection{Concept expansion}
A single BigQuery pull seeds two standard SNOMED concepts --- \textbf{318800} (GERD) and
\textbf{4223293} (GERD with esophagitis) --- and expands to all standard, selectable descendants via
the \code{cb_criteria} hierarchy path match, restricted to Fitbit-sharers. Because the pull is already
GERD-only, \code{has_gerd} is computed over \emph{all} returned rows: re-filtering on the two seed IDs
would discard every descendant-coded case and undercount GERD. The oesophagitis subset is identified
within the same pull by a concept-name match, which is robust to release-to-release changes in
descendant IDs.

\subsection{Case definitions}
\begin{table}[h]\centering\footnotesize
\begin{tabular}{@{}p{3.4cm}p{4.4cm}p{5.6cm}@{}}
\toprule
Variable & Rule & Role \\ \midrule
\code{has_gerd} / \code{has_esophagitis} & $\geq 2$ records, first $\geq 180$ days after the first
Fitbit record & \textbf{Primary} --- matches the published IBS paper's phenotype \\
\code{has_gerd_sens} & $\geq 2$ records at any time & Sensitivity (always retained) \\
\code{has_gerd_dates} & $\geq 2$ \emph{distinct} dates & Specificity check \\
\code{severe_gerd_binary} & GERD \emph{and} $\geq 2$ post-diagnosis PPI/$H_2$RA dates & Optional
treated-disease sensitivity \\
\bottomrule
\end{tabular}
\caption{Outcome variants. \code{GERD_PRIMARY_DEF} switches which rule is primary; the alternative is
always kept as \code{*_sens}.}
\end{table}

\noindent The temporal rule is the primary phenotype because it is what the published IBS paper used,
and because it is the only available guard against the reverse-causal edge (reflux $\rightarrow$
disturbed sleep) that covariate adjustment cannot address.

%======================================================================
\section{Cohort, Exposures, and Covariates}
%======================================================================
\subsection{Eligibility thresholds}
\begin{table}[h]\centering\footnotesize
\begin{tabular}{@{}p{4.2cm}p{2.6cm}p{6.6cm}@{}}
\toprule
Rule & Threshold & Rationale \\ \midrule
Main sleep only & \code{is_main_sleep} & Naps have different architecture \\
Valid sleep minutes & $(0, 1440]$ & Excludes device/sync errors \\
Short-sleep exclusion & $<30\%$ nights $<$4h & Removes wear artifacts without discarding occasional short nights \\
Sleep coverage & $\geq 180$ nights & $\sim$6 months stabilises nightly means \\
Valid activity day & $\geq 10$ wear h; 100--45,000 steps & Excludes non-wear and sensor spikes \\
Activity coverage & $\geq 30$ valid days & Balances stability against sample size \\
Comorbidity rule & $\geq 2$ distinct pre-Fitbit dates & Enforces temporality and specificity \\
Medication rule & $\geq 2$ dates in prior 12 months & Sustained rather than incidental use \\
\bottomrule
\end{tabular}
\caption{Thresholds, inherited from the antecedent studies for comparability.}
\end{table}

\subsection{Exposures}
Sleep metrics (minutes asleep, in bed, awake, restless, deep, light, REM, plus efficiency
$=$ asleep$/$in-bed) and activity metrics (daily steps, lightly/fairly/very/total active minutes,
sedentary minutes, corrected maximum daily heart rate) are averaged within participant and then
discretised into cohort quartiles with \textbf{Q1 as reference}. Quartiles avoid a linearity
assumption, resist outliers, and reveal dose--response shape. Each exposure is modelled
\emph{separately} because exposures within a domain are collinear.

\subsection{Combined design}
The combined analysis uses the \emph{intersection} cohort ($\geq 180$ sleep nights \emph{and}
$\geq 30$ activity days). It fits (A) each of the 15 exposures separately, and (B)
mutually-adjusted models pairing one sleep with one activity exposure, so each domain is adjusted for
the other:
\[
\logit(p) = \beta_0 + \sum_{j=2}^{4}\beta^S_j \mathbb{1}[S{=}Q_j]
+ \sum_{j=2}^{4}\beta^A_j \mathbb{1}[A{=}Q_j] + \sum_k \gamma_k Z_k .
\]
Four pairs are examined: minutes-asleep$\times$steps, efficiency$\times$very-active,
deep-sleep$\times$total-active, REM$\times$sedentary. Covariate timing for this cohort is anchored to
the sleep stream (always defined for its members).

\subsection{Covariate strategy}
\label{sec:covstrategy}
The adjustment set matches the published IBS paper --- age category, sex at birth, race, ethnicity,
education, income, alcohol, smoking, median BMI, depression, anxiety, Charlson index --- \emph{plus}
two GERD-specific gastrointestinal covariates, peptic ulcer disease (\code{has_pud}) and irritable
bowel syndrome (\code{has_ibs}), both clinically entangled with reflux. Acid-suppression therapy is
deliberately excluded from adjustment as a probable \emph{mediator}; it is used only to define the
treated-GERD sensitivity outcome.

\textbf{Pre-built covariates are reused wherever they exist.} This is a deliberate methodological
choice, not merely an optimisation. The covariate derivations are outcome-agnostic and already
underpin the published IBS paper; reusing them guarantees the GERD analysis is directly comparable to
it, rather than differing through incidental re-derivation. It also avoids re-reading multi-gigabyte
tables --- \code{R9_drug_df} alone is $\approx$3.1 GB in memory, and loading the sleep, activity,
steps and drug tables together approaches 9 GB. When a pre-built file is absent the covariate is
derived from the raw tables instead, and the code reports which path it took.

%======================================================================
\section{Software Architecture}
%======================================================================
The implementation is layered so that all six configurations share one statistical engine and one
data-loading path; only the cohort, exposure list and outcome column differ between them.

\begin{table}[h]\centering\footnotesize
\begin{tabular}{@{}p{5.2cm}p{8.4cm}@{}}
\toprule
File & Role \\ \midrule
\code{GERD Outcome Data Upload R9.R} & BigQuery pull of GERD records (both phenotypes) $\rightarrow$
\code{R9_gerd_outcome.rds}. Run once. \\
\code{GERD Data Prep R9.R} & \textbf{Data layer.} Schema-aware loading of cohort, exposures and the
one-row-per-participant covariate bundle; prefer-pre-built-else-derive; NULL-safe joins; memory
release after large tables are reduced. \\
\code{GERD Analysis Helpers R9.R} & \textbf{Engine.} Outcome definitions, exposure lists, adjustment
set, labels; descriptive/univariable/multivariable routines; sparsity guard; diagnostics; the
manuscript tables and figures. \\
\code{Sleep and GERD Analysis File R9.R} & Driver: C1 + C2 \\
\code{Activity and GERD Analysis File R9.R} & Driver: C3 + C4 \\
\code{Activity Sleep and GERD Analysis File R9.R} & Driver: C5 + C6, plus mutually-adjusted pairs \\
\code{GERD Pharmacologics Data Upload.R} & Optional PPI (21600095) / $H_2$RA (21600096) pull \\
\bottomrule
\end{tabular}
\caption{Code layers. Each driver sources the data layer and the engine, then calls one function.}
\end{table}

\begin{lstlisting}[caption={A driver reduces to this.}]
source("GERD Analysis Helpers R9.R")     # engine
source("GERD Data Prep R9.R")            # data layer

R9_gerd_outcome        <- read_first_existing("R9_gerd_outcome.rds", required = TRUE)
R9_gerd_outcome_status <- build_gerd_outcomes(R9_gerd_outcome,
                            first_fitbit_sleep_date_df, "first_fitbit_sleep_date")
covars_df              <- gerd_load_covariates(anchor = "sleep", ...)
# ... assemble one row per participant, create integer quartiles ...
modeling_df   <- prep_modeling_df(final_analysis_sleep_gerd_df, SLEEP_EXPOSURES)
sleep_results <- analyze_all_outcomes(modeling_df, exposures = SLEEP_EXPOSURES,
                                      coverage_var = "n_valid_nights", stub = "sleep")
\end{lstlisting}

\noindent Configuration lives in one place, at the top of the engine:
\code{GERD_PRIMARY_DEF} (default \code{"post_fitbit"}), \code{GERD_POST_FITBIT_LAG_DAYS} (180),
\code{GERD_P_ADJUST} (\code{"bonferroni"}), \code{GERD_MIN_CELL} (1), \code{GERD_OUTPUT_DIR}.

%======================================================================
\section{Statistical Analysis}
%======================================================================
\subsection{Models}
Descriptive tables are stratified by each outcome, reporting median [IQR] for continuous variables and
$n$ (\%) for categorical, with $t$-tests and chi-square tests as in the antecedent paper. Univariable
logistic regressions give unadjusted odds ratios for every exposure and covariate. For each exposure a
separate adjusted model is fit,
\[
\logit\!\big(\Pr[Y=1]\big) = \beta_0 + \sum_{j=2}^{4}\beta_j \mathbb{1}[X{=}Q_j] + \sum_k \gamma_k Z_k ,
\]
so $e^{\beta_j}$ is the adjusted odds ratio for quartile $j$ versus Q1. Reporting all three contrasts
preserves the \emph{shape} of the relationship rather than compressing it to a slope. Alongside each
full model a backward-AIC model is fit with the exposure forced to remain
(\code{MASS::stepAIC}), as a parsimony check: a stable estimate across both models indicates the
association is not an artifact of one covariate set. Each model reports complete-case $N$, ROC AUC
(\code{pROC}) and generalised VIF (\code{car}). Multiplicity is controlled by Bonferroni.

\subsection{Separation and sparsity guard}
\label{sec:guard}
Two features of this cohort make naive fitting dangerous: \code{has_pud} is recorded for only
\textbf{137 participants} cohort-wide, and the oesophagitis outcome is itself uncommon. A covariate
level containing zero events yields an infinite, meaningless odds ratio and destabilises the whole
model.

The engine therefore screens every model before fitting. A covariate is dropped \emph{from that model
only} if it is constant, single-level, or has any empty outcome$\times$level cell; each drop is
reported and recorded on the result object. This converts a silent failure into an explicit,
auditable decision. In validation the guard fired exactly as intended on the rare-outcome models,
dropping \code{has_pud} and several sparse demographic factors. \textbf{If it fires on the primary
models with real data, it must be reported in the manuscript's methods.}

\begin{lstlisting}[caption={The guard, applied inside the per-exposure fitting function.}]
guard      <- drop_unstable_covariates(d, outcome_col, covars, context = ...)
covars_use <- guard$keep       # constant / single-level / empty-cell terms removed
m_full     <- glm(as.formula(paste0(outcome_col, " ~ ",
                    paste(c(exposure, covars_use), collapse = " + "))),
                  data = d, family = binomial())
\end{lstlisting}

\subsection{Diagnostics}
Five assumption checks run for every configuration and are returned as numbers for the Results prose:
independence (duplicate identifiers), multicollinearity (adjusted GVIF range), separation (large
standard errors), influential points (maximum Cook's distance vs $4/N$), and linearity in the logit
(Box--Tidwell for median BMI and continuous age).

\subsection{Missing data and multiplicity}
Binary EHR indicators are coalesced \code{NA}$\rightarrow$\code{FALSE}: in EHR phenotyping the absence
of a qualifying record means the participant does not meet the phenotype. Survey and measurement
covariates are \emph{not} imputed; an explicit ``Unknown/Missing'' level is carried through and
reported, so differential missingness stays visible. Each adjusted model uses complete cases over its
own variables, so the effective $N$ is reported per model.

%======================================================================
\section{Manuscript Outputs}
%======================================================================
Running any driver writes, for \emph{both} outcomes, into \code{./manuscript_output/}:

\begin{table}[h]\centering\footnotesize
\begin{tabular}{@{}p{6.2cm}p{7.4cm}@{}}
\toprule
File & Manuscript element \\ \midrule
\code{*_table1.csv} & \textbf{Table 1} --- demographics by outcome group, median [IQR] / $n$ (\%) \\
\code{*_table2.csv} & \textbf{Table 2} --- exposure averages interleaved with quartiles, cutoffs
printed in the row labels \\
\code{*_supp_table1_univariate.csv} & \textbf{Supplement Table 1} --- unadjusted ORs \\
\code{*_supp_table2_multivariable.csv} & \textbf{Supplement Table 2} --- adjusted ORs, exposure rows,
Q1 shown as ``--'' \\
\code{*_figure1_forest.png} & \textbf{Figure 1} --- forest plots (Q1 reference) \\
\code{*_figure2_gvif.png} & \textbf{Figure 2} --- GVIF multicollinearity chart \\
\code{*_diagnostics.csv} & GVIF range, max Cook's distance, Box--Tidwell $p$, AUC range \\
\code{*_quartile_cutoffs.csv} & Exact unrounded boundaries for the Table 2 footnote \\
\code{combined_*_mutually_adjusted_pairs.csv} & Sleep$\times$activity pair estimates (C5/C6) \\
\bottomrule
\end{tabular}
\caption{Prefixes are \code{sleep_}, \code{activity_}, \code{combined_}, each followed by \code{gerd_}
or \code{esophagitis_}.}
\end{table}

Table 2 is emitted in the target format, with quartile boundaries computed from the data and written
into the row labels:

\begin{lstlisting}[caption={Table 2 excerpt as produced by the code.}]
Number of valid nights                     886 [532, 1,243]
Average sleeping minutes                     398 [369, 430]
Minutes asleep quartiles (minutes)
    Q1 (<369)                                    702 (25%)
    Q2 (369-398)                                 705 (25%)
    Q3 (399-430)                                 692 (25%)
    Q4 (>430)                                    693 (25%)
Daily steps quartiles (steps)
    Q1 (<5,785)   Q2 (5,785-7,489)   Q3 (7,490-9,108)   Q4 (>9,108)
\end{lstlisting}

%======================================================================
\section{Verification}
\label{sec:verify}
%======================================================================
The analysis code was validated in three stages before release.

\textbf{1. Schema validation.} Every dataset the code touches was opened and its columns and types
confirmed, rather than assumed --- covering the raw Fitbit, EHR, survey and person tables and all
pre-built covariate objects.

\textbf{2. Schema-exact synthetic data.} Synthetic datasets were constructed with column names and
types checked field-by-field against the real schema templates, including the deliberate hazards: a
duplicated \code{n_valid_days} column, the \code{answer_concept_id.x/.y} pair, en-dashed factor
levels, and an ultra-rare \code{has_pud} (20 of 3,000, mirroring the real 137).

\textbf{3. End-to-end execution.} All three drivers were run to completion on that data. Final status:
\textbf{exit code 0 for all three}, 50 output files, zero errors and zero warnings.

\subsection{Defects found and fixed}
Five defects surfaced during schema alignment. Each would have broken the run or silently corrupted
results:

\begin{table}[h]\centering\footnotesize
\begin{tabular}{@{}p{4.6cm}p{9.0cm}@{}}
\toprule
Defect & Consequence and fix \\ \midrule
\code{library(MASS)} masked \code{dplyr::select()} & Every \code{select()} downstream failed
immediately. MASS/car/pROC are now referenced with \code{::} and never attached; dplyr attaches last. \\
Duplicate \code{n_valid_days} & Joining the wear-hours and steps summaries produced
\code{n_valid_days.x/.y}, breaking the coverage column. The wear file is now slimmed to two columns. \\
\code{answer_concept_id.x/.y} & Polluted the covariate bundle; only \code{alcohol_likert_final} is taken. \\
Education/income from free text & The \code{*_response} strings did not match the antecedent script's
matching, silently collapsing every participant to ``Missing''. Categories are now mapped from
\emph{concept IDs}. \\
En-dashed factor levels & \code{cci_cat} levels failed to match and emitted 50+ warnings. Dash style is
now normalised. \\
\bottomrule
\end{tabular}
\caption{Defects found by schema alignment and end-to-end execution.}
\end{table}

\subsection{Scope of the verification}
This establishes that the code runs correctly against the real \emph{schema} and emits the required
format. It does not, and cannot, anticipate the missingness patterns and case counts of the real data.
Two console outputs should be checked on the first production run: the eligible sleep cohort should
print $N = 19{,}995$ (the published cohort; a different value means a different
\code{valid_population_sleep} was picked up), and \code{build_gerd_outcomes()} prints case counts under
both definitions --- if the post-Fitbit count is very small, the ``ever'' rule may be more appropriate.

%======================================================================
\section{Sensitivity Analyses and Execution}
%======================================================================
\begin{enumerate}[leftmargin=1.4em,itemsep=1pt]
  \item \textbf{Alternative case definition} (\code{has_gerd_sens}): the ``$\geq 2$ records ever'' rule,
  always retained alongside the primary temporal rule.
  \item \textbf{Distinct-date rule}: $\geq 2$ records on distinct calendar dates, guarding against
  same-day code inflation.
  \item \textbf{Sleep$\times$activity overlap}: analyses restricted to concurrently well-measured days.
  \item \textbf{Treated (``severe'') GERD}: GERD \emph{and} $\geq 2$ PPI/$H_2$RA prescription dates
  after the first diagnosis --- an EHR proxy for clinically managed disease, mirroring the IBS study's
  severe-IBS construction. Guarded by file existence, so the primary pipeline runs either way.
\end{enumerate}

\subsection{Execution order}
\begin{enumerate}[leftmargin=1.4em,itemsep=0pt]
  \item Verify the GERD seed concepts in the Cohort Builder.
  \item \code{source("GERD Outcome Data Upload R9.R")} $\rightarrow$ \code{R9_gerd_outcome.rds} (once).
  \item Run any of the three drivers; each produces both outcomes.
  \item Optionally run the pharmacologics pull, then re-run the sleep driver for the treated-GERD
  sensitivity.
\end{enumerate}
All seven \code{.R} files and the \code{.rds} datasets must share one working directory. Detailed
instructions, settings and troubleshooting are in \code{HOW_TO_RUN_GERD_ANALYSIS.md}.

%======================================================================
\section{Interpretation Framework}
%======================================================================
An adjusted OR of, say, $0.85$ for Q4 versus Q1 of \code{min_asleep_quartile} means participants in
the highest sleep-duration quartile have $15\%$ lower adjusted odds of the outcome. A monotone decline
across Q2$\rightarrow$Q3$\rightarrow$Q4 is a dose--response signature.

\begin{table}[h]\centering\footnotesize
\begin{tabular}{@{}p{4.4cm}p{2.4cm}p{6.4cm}@{}}
\toprule
Exposure (high vs low) & Predicted aOR & Primary mechanism \\ \midrule
More minutes asleep & $<1$ & Clearance, less arousal, lower sensitivity \\
Higher sleep efficiency & $<1$ & Consolidated sleep, less low-clearance recumbency \\
More deep / REM sleep & $<1$ & Restorative sleep, autonomic balance \\
More minutes awake / restless & $>1$ & Fragmentation, sympathetic arousal \\
More steps / active minutes & $<1$ & Weight control, vagal tone, motility \\
More sedentary minutes & $>1$ & Adiposity, postprandial recumbency \\
Higher corrected max HR & $<1$ & Cardiovascular fitness proxy \\
\bottomrule
\end{tabular}
\caption{Pre-specified directional predictions, applying to both outcomes.}
\end{table}

Three consistency checks structure interpretation. \emph{Across outcomes}: an association holding for
both broad GERD and the erosive subset is less likely to be an artifact of symptom reporting; one
confined to symptom-coded GERD points to a perceptual rather than structural pathway. \emph{Across
domains}: agreement between sleep-only, activity-only and mutually-adjusted estimates indicates
neither domain is merely a proxy for the other. \emph{Across case definitions}: agreement between the
primary temporal rule and the alternatives indicates robustness to phenotype choice. Because the
design is cross-sectional, estimates remain associations, not causal effects.

%======================================================================
\section{Limitations}
%======================================================================
\begin{itemize}[leftmargin=1.3em,itemsep=1pt]
  \item \textbf{Cross-sectional design.} The temporal phenotype mitigates but does not eliminate
  reverse causation; reflux itself disrupts sleep.
  \item \textbf{EHR phenotyping.} Diagnosis codes miss undiagnosed and over-the-counter--managed
  reflux and vary with coding practice; the $\geq 2$-record rule trades sensitivity for specificity.
  Participants with more health-care contact are both more likely to be coded and may differ in sleep
  and activity --- partly blocked by adjusting for comorbidity and socio-economic status.
  \item \textbf{Name-based erosive definition.} Identifying oesophagitis by concept name may miss
  erosive concepts coded outside the GERD hierarchy unless their seeds are added.
  \item \textbf{Selection.} The cohort owns a Fitbit, shares it and the EHR, and wears the device
  enough to meet coverage thresholds --- younger, more affluent and more health-engaged than the
  general population.
  \item \textbf{Residual confounding.} Diet, meal timing, caffeine, body position and hiatal hernia
  are unmeasured.
  \item \textbf{Precision.} Oesophagitis is rarer than broad GERD and the combined cohort is the
  smallest, so those configurations are the most precision-limited; the sparsity guard
  (Section~\ref{sec:guard}) makes any resulting instability explicit rather than silent.
  \item \textbf{Covariate timing.} \code{has_pud} derives from a file computed on the activity-based
  anchor, a small documented inconsistency for the sleep analysis.
\end{itemize}

%======================================================================
\section{Conclusion}
%======================================================================
This specification defines a reproducible suite relating objectively measured sleep and physical
activity to two nested EHR-defined outcomes in \emph{All of Us} CDR v9. Six configurations are
executed by a single engine over a single schema-verified data layer, so differences between them
reflect biology rather than analytic drift. The code is written against verified schemas, reuses the
covariate objects underpinning the published IBS paper for direct comparability, guards explicitly
against separation in the rare-outcome models, and emits every table and figure in the target
manuscript's format. Together with the temporal, distinct-date, overlap and treated-disease
sensitivity analyses, the design supports a robust assessment of whether sleep and activity are
associated with reflux disease and its erosive form.

\appendix
%======================================================================
\section{Appendix: Concept Reference}
%======================================================================
\begin{table}[h]\centering\footnotesize
\begin{tabular}{@{}p{5.2cm} p{4.6cm} p{3.4cm}@{}}
\toprule
Purpose & Concept ID(s) & Note \\ \midrule
GERD (outcome seeds) & 318800; 4223293 & Descendant-expanded \\
Oesophagitis subset & 4223293 + name match & Erosive subset \\
BMI measurement & 3038553 & Body mass index \\
IBS (covariate) & 75576; 4234788; 4261072; 4057826 & General + subtypes \\
Peptic ulcer disease & 4318534; 4134146; 4265600; 4027663 & Via \code{R9_pud_status} \\
Depression / anxiety & 440383 / 441542 & \\
Diabetes / hypertension & 201820 / 316866 & \\
Heart failure / MI / stroke & 316139 / 4329847 / 381316 & \\
COPD / sleep apnea & 255573 / 313459 & \\
PPI / $H_2$RA & 21600095 / 21600096 & Drug-class ancestors \\
Beta-blockers / CCBs & 21601664 / 21601765 & \\
Stimulants / antidepressants & 21604753 / 21604686 & \\
Antipsychotics & 21604490 & \\
Anxiolytics & 21604600; 21604565 & \\
Hypnotics & 21604635; 21604653; 21604685; 21604661 & \\
\bottomrule
\end{tabular}
\end{table}

%======================================================================
\section{Appendix: Key Derived Variables}
%======================================================================
\begin{longtable}{@{}p{4.8cm}p{8.8cm}@{}}
\toprule
Variable & Meaning \\ \midrule \endhead
\code{first_fitbit_sleep_date} / \code{first_fitbit_date} & Anchor for the temporal outcome rule and all pre-Fitbit covariate windows \\
\code{n_valid_nights} / \code{n_valid_days} & Coverage counts driving eligibility \\
\code{avg_min_*}, \code{avg_sleep_efficiency} & Participant-level sleep means \\
\code{avg_daily_steps}, \code{avg_*_active_min}, \code{avg_sedentary_min} & Participant-level activity means \\
\code{*_quartile} & Cohort quartile of the corresponding metric (Q1 reference) --- the modelled exposures \\
\code{has_gerd} / \code{has_esophagitis} & Primary outcomes under \code{GERD_PRIMARY_DEF} \\
\code{has_gerd_sens} & The alternative case definition, always retained \\
\code{severe_gerd_binary} & GERD and $\geq 2$ post-diagnosis PPI/$H_2$RA dates \\
\code{has_pud} / \code{has_ibs} & GERD-specific gastrointestinal covariates \\
\code{cci_score} / \code{cci_cat} & Charlson index, continuous and categorical \\
\code{median_bmi} & Within-person median BMI \\
\code{education_collapsed} / \code{income_collapsed} & Socio-economic factors mapped from concept IDs \\
\code{covars_dropped} & Covariates removed from a model by the sparsity guard \\
\code{R9_final_analysis_*_gerd_df} & The three one-row-per-participant modelling frames \\
\bottomrule
\end{longtable}

\vspace{0.6em}
\noindent\rule{\linewidth}{0.4pt}\\
{\footnotesize \emph{End of report.} This document describes code and methodology only; it contains no
individual-level \emph{All of Us} data.}

\end{document}
